% Level 2, Week 1 -- Welcome Back, Python
\lesson{1}{Welcome Back, Python}{Every value has a type, and f-strings glue values into friendly messages.}
\lessoninfo{Concepts: data types (I), f-strings (I), input (R), if/elif/else (R), and/or/not (I)}{Tools: Thonny, value cards}{CSTA: 2-AP-11, 2-AP-12, 2-NI-05}

\begin{objectives}
  \item The student can sort values into the four types \texttt{int}, \texttt{float}, \texttt{str} and \texttt{bool} and check a guess with \texttt{type()}. (L2.G1; L2.W01.1)
  \item The student can build a message with an f-string and turn typed text into a number with \texttt{int()}. (L2.G1; L2.W01.2)
  \item The student can combine conditions with \texttt{and}, \texttt{or} and \texttt{not} inside \texttt{if}/\texttt{elif}/\texttt{else}. (L2.G2; L2.W01.3)
  \item The student can read a \texttt{TypeError} from the bottom up and explain why \texttt{"5" + 5} fails. (L2.G6; L2.W01.4)
\end{objectives}

\begin{materials}
  \item Print (one per student): the \textbf{Placement quiz in the Assessments booklet}, workbook chapter 1, homework 1.
  \item Value cards (5 minutes to make): write each of these on its own sticky note or index card:
        \texttt{7}\quad \texttt{3.5}\quad \texttt{"cat"}\quad \texttt{True}\quad \texttt{"7"}\quad \texttt{-2}\quad \texttt{0.0}\quad \texttt{False}\quad
        \texttt{"True"}\quad \texttt{100}\quad \texttt{"3.5"}\quad \texttt{2.0}\quad \texttt{0}\quad \texttt{" "}.
        Then four bin labels: \textbf{int} (whole number), \textbf{float} (decimal number), \textbf{str} (text in quotes), \textbf{bool} (True or False).
  \item On each Mac: the \texttt{kids-coding} folder from Level 1. In Thonny open
        \texttt{2-intermediate/exercises/week01/ex-hero-card/starter.py} (\keys{Cmd}+\keys{O}). Run \texttt{check\_my\_work.py} once; if it asks
        for pytest: \menu{Tools} → \menu{Manage packages…} → \texttt{pytest} → \menu{Install} (no account needed).
  \item A timer, paper and pencils.
\end{materials}

\needspace{16\baselineskip}
\section{At a glance}
Times are a guide: set a timer for each phase and move on when it rings.
\begin{timeplan}
  \timeslot{0--15}{Warm-up}{Welcome back (2 min), then the 15-minute Placement quiz from the Assessments booklet (pencil and paper, no computers). Collect it; score it after the session.}
  \timeslot{15--23}{Concept}{Unplugged \enquote{Type sorting} with the value cards. Name the four types. Personal-information talk: hero names, not real names.}
  \timeslot{23--31}{Guided demo}{In the Thonny Shell: \texttt{type()}, the \texttt{"5" + 5} bug, \texttt{int()}, f-strings. Then run the Hero Card starter together.}
  \timeslot{31--53}{Hands-on build}{Pairs: Exercise 1.1 Hero Card, then 1.2 Mood Meter. Swap driver and navigator every 8 minutes. Last 5 minutes: free-build.}
  \timeslot{53--60}{Wrap-up}{Show one hero card each, exit ticket, reflection, hand out homework 1.}
\end{timeplan}

\begin{teachernote}[About the placement quiz]
It checks what stuck from Level 1 and is not graded: say so. Answers and scoring guide: Answer Key.
Use the result to pick the support or extension path below. You may instead give it earlier in the week and start here at the Concept phase.
\end{teachernote}

\section{Warm-up (0--15 min)}
\begin{teacherscript}[Say]
\enquote{Welcome to Level 2! You already know a lot of Python: loops, variables, questions with \texttt{input},
choices with \texttt{if}, and your own commands with \texttt{def}. Today we start with a short quiz so I know
what to help with. It is not graded. If you do not know an answer, write \enquote{not sure} and move on. That is
useful to me too.}
\end{teacherscript}
Hand out the quiz, start a 15-minute timer, and sit with the group. Do not help with answers; you may read a question aloud.

\section{Concept (15--23 min)}
\begin{unplugged}[Type sorting]
Lay the four bin labels on the table. Shuffle the value cards and deal them out. Students take turns placing a
card in a bin and saying why. Others give a thumbs up or a \enquote{wait!}.
\par\textbf{Answers.} int: \texttt{7}, \texttt{-2}, \texttt{100}, \texttt{0}.\quad float: \texttt{3.5}, \texttt{0.0}, \texttt{2.0}.\quad
bool: \texttt{True}, \texttt{False}.\quad str: \texttt{"cat"}, \texttt{"7"}, \texttt{"True"}, \texttt{"3.5"}, \texttt{" "}.
\par\textbf{The trick cards} are \texttt{"7"}, \texttt{"True"} and \texttt{"3.5"}: the quotes make them text, even though they look like a number or a bool.
\texttt{" "} (a space in quotes) is text too, and so is \texttt{""} (empty text).
\end{unplugged}

\begin{teacherscript}[Say]
\enquote{Every value in Python has a \textbf{type}. The type tells Python what it is allowed to do with the value.
You can add two numbers. You can glue two pieces of text together. But Python gets confused if you mix them up.
Quotes are the big clue: anything inside quotes is text, which Python calls a \textbf{string}, or \texttt{str} for short.}
\end{teacherscript}

\begin{analogy}[Bricks and stickers]
A number like \texttt{5} is like a LEGO brick: you can click bricks together to make a bigger one (\texttt{5 + 5} is \texttt{10}).
Text like \texttt{"5"} is a sticker with a 5 printed on it. You can put stickers next to each other
(\texttt{"5" + "5"} is \texttt{"55"}), but you cannot click a sticker onto a brick. That is exactly the error Python shows for \texttt{"5" + 5}.
\end{analogy}

\begin{safety}[Hero names, not real names]
Ask: \enquote{Which of these would be OK to type into a game we make?} Read out: a made-up hero name, your favorite color,
your full real name, the name of your school, your home address, your birthday, your favorite animal.
Agree together: made-up names, colors and favorites are fine; real names, schools, addresses and birthdays are
\textbf{personal information}, which we keep private. All our programs ask for a \enquote{hero name} or \enquote{nickname} instead.
Tell students: if any app or website asks for personal information, they check with you first.
\end{safety}

\section{Guided demo (23--31 min)}
Project or share one screen. Type into the \textbf{Shell} (the lower part of Thonny, after the \texttt{>>>} prompt) and press \keys{Return} after each line.
Ask the students to predict each answer \emph{before} you press \keys{Return}.

\begin{workedexample}[Types, a bug and an f-string in the Shell]
\begin{pycode}[Type these one at a time into the Shell]
type(7)              # <class 'int'>
type("7")            # <class 'str'>
type(3.5)            # <class 'float'>
type(True)           # <class 'bool'>
"5" + 5              # TypeError! (read it from the bottom line up)
"5" + "5"            # '55'  stickers next to each other
int("5") + 5         # 10    int() turns the text into a number
name = "Nova"
level = 4
f"{name} is level {level}"     # 'Nova is level 4'
level >= 3 and name == "Nova"  # True
\end{pycode}
\begin{enumerate}[leftmargin=1.4em]
  \item After the \texttt{TypeError}, point at the \textbf{last line} of the red text: \enquote{This line says WHAT went wrong.
        \texttt{can only concatenate str (not "int") to str} means: I can only glue text to text.}
  \item For the f-string, say: \enquote{The \texttt{f} in front means \enquote{fill in}. Python fills each pair of curly brackets
        with the value of the variable inside them.}
  \item Open \texttt{ex-hero-card/starter.py}, press \menu{Run} (the green button, or \keys{F5}), answer the three questions and
        watch it crash. Read the traceback together from the bottom up: the error, then the line number. Do not fix it yet:
        that is the students' first job.
\end{enumerate}
\end{workedexample}
The finished Hero Card, for you (students build it themselves):
\pyfile{../exercises/week01/ex-hero-card/solution.py}

\section{Hands-on build (31--53 min)}
\begin{pairroles}[Driver and navigator]
The \textbf{driver} types. The \textbf{navigator} reads the TODO, spots typos and says \enquote{run it now!} after every small change.
Swap when the 8-minute timer rings (at about minute 39 and 47).
\textbf{1 student:} you are the navigator; ask questions, do not type.
\textbf{2 students:} one pair. \textbf{3 students:} a rotating trio: driver, navigator and \textbf{tester} (the tester chooses
the answers to type when the program runs, including tricky ones like level 10).
\textbf{4 students:} two pairs; at the end each pair runs the other pair's program.
\end{pairroles}

\begin{enumerate}[leftmargin=1.4em]
  \item \textbf{Exercise 1.1 Hero Card} (\texttt{week01/ex-hero-card}, about 12 minutes). Fix the \texttt{TypeError} (TODO 1), then add the f-string lines.
  \item \textbf{Exercise 1.2 Mood Meter} (\texttt{week01/ex-mood-meter}, about 10 minutes). Students fill in the prediction table in the workbook before running.
  \item After each exercise: open \texttt{check\_my\_work.py} in the same folder and press \menu{Run}. Green \texttt{[PASS]} lines mean done.
\end{enumerate}

\begin{bughunt}[\texttt{"5" + 5}]
The Hero Card starter crashes on the line \texttt{print("  Next level:   " + level + 1)} because \texttt{input()} always gives back text.
Ask: \enquote{What type is \texttt{level}? How do we know?} Let students add \texttt{print(type(level))} just above the crash line to find out.
Two good fixes: wrap the input in \texttt{int(...)} (TODO 1), and use an f-string for the line: \texttt{f"  Next level:   \{level + 1\}"}.
\end{bughunt}

\begin{freebuild}[Make the card yours]
Last 5 minutes: add a sidekick name, a favorite color, a border made with \texttt{"+" + "-" * 30 + "+"}, or a fourth mood
(\enquote{Sleepy} for energy 0 or 1) in the Mood Meter. Everything stays made-up: hero names only.
\end{freebuild}

\section{Wrap-up (53--60 min)}
Each student (or pair) shows one hero card and says one type they used.

\begin{formativecheck}[Exit ticket]
Write on a sticky note: what does this print, and how would you fix it?
\begin{pycode}
coins = input("Coins? ")     # the player types 3
print(coins + 1)
\end{pycode}
\textbf{Answer:} it crashes with a \texttt{TypeError}, because \texttt{coins} is the text \texttt{"3"}. Fix:
\texttt{coins = int(input("Coins? "))}. Then it prints \texttt{4}.
\textbf{Look for:} students who say \enquote{4} without spotting the type problem need the Concept cards again next week.
\end{formativecheck}

\begin{reflection}[Say it or write it]
\begin{itemize}[leftmargin=1.2em]
  \item Which value card surprised you? Why?
  \item What does the last line of a red error message tell you?
  \item When would you use \texttt{and}, and when \texttt{or}? Give a game example.
\end{itemize}
\end{reflection}
Hand out homework 1 (Dragon Door). It takes 15--20 minutes.

\section{Notes for the teacher}
\begin{commonmistakes}
  \item \textbf{Forgetting the \texttt{f}.} \texttt{"Hi \{name\}"} prints the curly brackets. Fix: \texttt{f"Hi \{name\}"}.
  \item \textbf{Using a number from \texttt{input()} without \texttt{int()}.} Text times a whole number repeats it (\texttt{"4" * 10} is
        \texttt{"4444444444"}) and text plus a number crashes.
  \item \textbf{\texttt{yes} without quotes.} \texttt{if answer == yes:} gives \texttt{NameError}. Text needs quotes: \texttt{"yes"}.
  \item \textbf{\texttt{=} instead of \texttt{==}.} One equals sign stores a value; two equals signs compare.
  \item \textbf{Half a condition.} \texttt{if energy >= 5 or 8:} is always true. Each side of \texttt{and}/\texttt{or} needs a full comparison:
        \texttt{energy >= 5 or energy == 8}.
  \item \textbf{\texttt{"True"} versus \texttt{True}.} The first is text (always counts as true in an \texttt{if}); the second is the bool.
\end{commonmistakes}

\begin{differentiation}
\supportpath Use the placement quiz: if loops or variables were shaky, sit with that pair for TODO 1 and 2 of the Hero Card and type the first f-string
together. Skip the Mood Meter in class; it can be a homework extra. Keep the value cards on the table as a reminder.
\extensionpath Try \texttt{type()} on \texttt{3 / 2}, \texttt{4 / 2} and \texttt{7 // 2} and explain the answers. Show health with one decimal place:
\texttt{f"\{health:.1f\}"}. In the Mood Meter, add a third question and a condition that uses \texttt{and}, \texttt{or} and \texttt{not} together.
\end{differentiation}

\begin{stretchchallenge}[Rank badges]
Add a badge to the hero card: Bronze for levels 1--4, Silver for 5--9, Gold for 10 and up. Use \texttt{if}/\texttt{elif}/\texttt{else} and
\texttt{and}, for example \texttt{level >= 5 and level <= 9}. Python also understands \texttt{5 <= level <= 9}: can they find out if it gives the same answer?
\end{stretchchallenge}

\begin{teachernote}[If you are short on time]
Cut the Mood Meter (it becomes an optional homework extra) and shorten Type sorting to the three trick cards plus one card per bin.
Keep the \texttt{"5" + 5} bug hunt and the exit ticket: they are the heart of this lesson.
\end{teachernote}
